% !PRINTED 732
\weberpage{732}
%<<CONT
gewissermaassen wie ein Fundamentalsystem verhalten, was durch
folgenden Satz bestimmter ausgedrückt wird:

2. \gesperrt{Wenn eine Normaleinheit des Körpers $H_m$
von der Form}
\[
\mathfrak{E}(r) = \tau_0^{e_0}\tau_1^{e_1}\ldots\tau_{\nu'-1}^{e_{\nu'-1}},
\]
\gesperrt{in der die Exponenten $e_0, e_1, \ldots, e_{\nu'-1}$ ganze rationale
 Zahlen sind, mit allen ihren conjugirten
Werthen einerlei Vorzeichen hat, so müssen die
sämmtlichen $e_0, e_1, \ldots, e_{\nu'-1}$ gerade Zahlen sein,
und $\mathfrak{E}(r)$ ist das Quadrat einer Normaleinheit.}

Den Beweis können wir wieder durch vollständige Induction
führen, wollen aber zunächst den Ausdruck für $\tau_\beta$ etwas umformen
.

Die conjugirten Werthe zu der Einheit $\mathfrak{E}(r)$ erhalten wir
in der Form
\[
\mathfrak{E}(r_\beta) = \tau_\beta^{e_0}\tau_{\beta+1}^{e_1}\ldots\tau_{\beta+\nu'-1}^{e_{\nu'-1}},
\]
und diese Ausdrücke sollen also für alle Werthe von $\beta$ dasselbe
Vorzeichen haben.

Es ist nach §.\,199, (5)
\[
\tau_\beta = \mathrm{tg}\,\frac{5^\beta\pi}{m} = \frac{1}{2}\frac{\sin\left(\dfrac{2\pi}{m}5^\beta\right)}{\left[\cos\left(\dfrac{5^\beta\pi}{m}\right)\right]^2}.
\]

Hierin setzen wir zur Abkürzung
\weq{14}{\sigma_\beta = \sin\left(\frac{2\pi}{m}5^\beta\right),}
so dass $\sigma_\beta$ immer dasselbe Vorzeichen hat, wie $\tau_\beta$.

Daraus bilden wir das Product
\weq{15}{S_\beta = \sigma_\beta^{e_0}\sigma_{\beta+1}^{e_1}\ldots\sigma_{\beta+\nu'-1}^{e_{\nu'-1}},}
und unsere Voraussetzung ist also die, dass $S_\beta$ für alle Werthe
von $\beta$ dasselbe Vorzeichen hat.

Es soll nachgewiesen werden, dass die ganzen Zahlen
$e_0, e_1, \ldots, e_{\nu'-1}$ alle gerade sein müssen. Der Satz ist
richtig für $m=8$; denn in diesem Falle ist
\[
\sigma_0 = \sin\frac{2\pi}{8}, \sigma_1 = \sin\frac{10\pi}{8} = -\sin\frac{6\pi}{8},
\]
%CONT>>
